% !TEX root = ../main.tex

\section{Construction of FEC Contribution Share Variables}
\label{sec:appendix_fecscreen}

This appendix provides a formal breakdown of the algorithm and filtering methodology used to construct the state-level political orientation metrics from the Federal Election Commission (FEC) campaign finance datasets. The screening framework eliminates intra-party conflict, non-individual contributions, conduit noise, and issue-based Political Action Committees (PACs) to accurately extract grassroots partisan leaning across presidential election years (1996, 2000, 2004, 2008, 2012, 2016, 2020).

\subsection{Candidate Filtering and Identification}
To isolate contributions that reflect general election partisan alignment, the candidate master universe is first restricted to the two major parties: Democrats (\texttt{DEM}) and Republicans (\texttt{REP}). For presidential cycles, inclusion is strictly limited to the ultimate primary winners who advanced to the general election, following \citet{meeuwis2022}. This prevents heightend intra-party competition (e.g., primary campaigns) from inflating a state's partisan polarization index. Voters supporting a losing primary candidate do not necessary support the general election nominee, and thus their contributions are not indicative of the state's overall partisan leaning in the general election.

Let $\mathcal{C}_t$ denote the set of valid presidential candidates in election cycle $t$:
\begin{align*}
    \mathcal{C}_{1996} &= \{\text{Bill Clinton (DEM)}, \text{Bob Dole (REP)}\} \\
    \mathcal{C}_{2000} &= \{\text{Al Gore (DEM)}, \text{George W. Bush (REP)}\} \\
    \mathcal{C}_{2004} &= \{\text{John Kerry (DEM)}, \text{George W. Bush (REP)}\} \\
    \mathcal{C}_{2008} &= \{\text{Barack Obama (DEM)}, \text{John McCain (REP)}\} \\
    \mathcal{C}_{2012} &= \{\text{Barack Obama (DEM)}, \text{Mitt Romney (REP)}\} \\
    \mathcal{C}_{2016} &= \{\text{Hillary Clinton (DEM)}, \text{Donald Trump (REP)}\} \\
    \mathcal{C}_{2020} &= \{\text{Joe Biden (DEM)}, \text{Donald Trump (REP)}\}
\end{align*}

\subsection{Multi-Stage Committee Screening Algorithm}
Committees ($c$) are classified into one of four categories: Democratic (\texttt{DEM}), Republican (\texttt{REP}), Excluded (\texttt{EXC}), or Indeterminate (\texttt{IND}). Classification follows a four-step sequence:

\subsubsection{Stage 1: Institutional Code Classification}
Initial filtering relies on the FEC committee type designations (\texttt{cmte\_tp}):
\begin{itemize}
    \item \texttt{EXC}: Communication costs (\texttt{C}), presidential primaries (\texttt{D}), electioneering communications (\texttt{E}), house committees (\texttt{H}), independent expenditors (\texttt{I}), and senate committees (\texttt{S}) are definitively removed.
    \item Committees categorized as party organizations (\texttt{X}, \texttt{Y}, \texttt{Z}) are automatically mapped to \texttt{DEM} or \texttt{REP} matching their reported party affiliation (\texttt{cmte\_pty\_affiliation}). If they report any other party affiliation, they are marked \texttt{EXC}.
    \item Standard PACs (\texttt{N}, \texttt{Q}), Super PACs (\texttt{O}), Hybrid PACs (\texttt{V}, \texttt{W}), single-candidate independent entities (\texttt{U}), and core presidential committees (\texttt{P}) are set as \texttt{IND} and sent to the next stage.
\end{itemize}

\subsubsection{Stage 2: Official Candidate Linkage}
For any committee still marked \texttt{IND}, if it possesses an official candidate link (\texttt{cand\_id}):
\begin{equation}
    \text{Status}(c) = \begin{cases} 
        \texttt{DEM} & \text{if } \texttt{cand\_id}_c \in \mathcal{C}_t \text{ and Party}(\texttt{cand\_id}_c) = \texttt{DEM} \\
        \texttt{REP} & \text{if } \texttt{cand\_id}_c \in \mathcal{C}_t \text{ and Party}(\texttt{cand\_id}_c) = \texttt{REP} \\
        \texttt{EXC} & \text{if } \texttt{cand\_id}_c \notin \mathcal{C}_t \text{ and } \texttt{cand\_id}_c \neq \emptyset \\
        \texttt{IND} & \text{otherwise} \\
    \end{cases}
\end{equation}

\subsubsection{Stage 3: Unofficial Linkage via Independent Expenditures}
For unlinked PACs, behavior is inferred using individual transaction data from the independent expenditure files (\texttt{pas2}). Valid expenditures supporting or opposing candidates are quantified. Let $N_{\text{for}, p}$ and $A_{\text{for}, p}$ represent the total transaction count and aggregate dollar volume spent \textit{supporting} party $p \in \{\text{DEM}, \text{REP}\}$, and let $N_{\text{agt}, p}$ and $A_{\text{agt}, p}$ represent expenditure amounts spent \textit{opposing} party $p$. 

We construct committee-level two directional indexes normalized to $[-1, 1]$, where $+1$ represents purely Republican support and $-1$ represents purely Democratic support:
\begin{equation}
    \mathcal{I}_{n}(c) = 2 \times \left( \frac{N_{\text{for, REP}} + N_{\text{agt, DEM}}}{N_{\text{total}}} \right) - 1
\end{equation}
\begin{equation}
    \mathcal{I}_{\text{sum}}(c) = 2 \times \left( \frac{A_{\text{for, REP}} + A_{\text{agt, DEM}}}{A_{\text{total}}} \right) - 1
\end{equation}
Committees with extreme, uniform alignment are captured as partisan vehicles:
\begin{equation}
    \text{Status}(c) = \begin{cases} 
        \texttt{DEM} & \text{if } \mathcal{I}_{n}(c) \le -0.75 \text{ and } \mathcal{I}_{\text{sum}}(c) \le -0.75 \\
        \texttt{REP} & \text{if } \mathcal{I}_{n}(c) \ge 0.75 \text{ and } \mathcal{I}_{\text{sum}}(c) \ge 0.75 
    \end{cases}
\end{equation}

\subsubsection{Stage 4: Text-Pattern Regex Targeted Capture}
Remaining \texttt{IND} committees are passed through a case-insensitive regular expression keyword check targeting party infrastructure and candidate names:
\begin{itemize}
    \item {Democratic Keywords}: Candidate last names in $\mathcal{C}_{t,\text{DEM}}$, \texttt{DNC}, \texttt{DSCC}, \texttt{DCCC}, \texttt{DEMOCRATIC}, \texttt{DEMOCRAT(S)}, \texttt{SENATE MAJORITY PAC}, \texttt{HOUSE MAJORITY PAC}.
    \item {Republican Keywords}: Candidate last names in $\mathcal{C}_{t,\text{REP}}$, \texttt{RNC}, \texttt{NRSC}, \texttt{NRCC}, \texttt{GOP}, \texttt{REPUBLICAN(S)}, \texttt{CONGRESSIONAL LEADERSHIP FUND}, \texttt{SENATE LEADERSHIP FUND}.
\end{itemize}

\subsubsection{Stage 5: Secondary Manual Screening and Overrides}
To guarantee that the selected active portfolio isolates core partisan alignment without introducing non-comparable baseline noise, a targeted manual screening is applied to high-volume committees:
\begin{itemize}
    \item Automated intermediary fundraising platforms (e.g., \texttt{ACTBLUE}) are explicitly set to \texttt{EXC} to eliminate baseline inflation, as individual allocations are already captured downstream at their final committee destinations.
    \item Ideological PACs that frequently back internal party primary candidates or prioritize single-issue advocacy (e.g., \texttt{EMILY'S LIST}, \texttt{FAIR FIGHT}, \texttt{AMERICANS FOR PROSPERITY ACTION}, \texttt{CLUB FOR GROWTH ACTION}) are systematically marked \texttt{EXC}.
    \item Industry-specific lobbying coalitions and profession-based bodies (e.g., automotive dealers, realtors, law associations, or corporate consulting PACs) are marked \texttt{EXC}..
    \item Specialized party or candidate-adjacent Super PACs with atypical names (e.g., \texttt{COMMITTEE FOR CHANGE}, \texttt{SMP}, \texttt{TAKE BACK THE HOUSE 2020}) are manually preserved as \texttt{DEM} or \texttt{REP} following manual verification.
\end{itemize}

\subsubsection{Final Set of committees}
\Cref{tab:B1} lists ten of the finally selected committees for each year, ordered by the number of total valid contributions. 

\begin{table}[tbp] \centering
\newcolumntype{R}{>{\raggedleft\arraybackslash}X}
\newcolumntype{L}{>{\raggedright\arraybackslash}X}
\newcolumntype{C}{>{\centering\arraybackslash}X}

\caption{Top 10 Selected Committees by Election Year}
\label{tab:B1}

\footnotesize
\begin{tabularx}{\linewidth}{c L c r r}

\toprule
{Year}&{Committee Name}&{Party}&{N}&{Amount} \tabularnewline
\midrule \addlinespace[\belowrulesep]
1996&REPUBLICAN NATIONAL COMMITTEE - RNC&REP&61957&39243617 \tabularnewline
&NATIONAL REPUBLICAN CONGRESSIONAL COMMITTEE CONTRIBUTIONS&REP&55651&21067394 \tabularnewline
&DOLE FOR PRESIDENT INC&REP&32673&19303022 \tabularnewline
&CLINTON/GORE '96 PRIMARY COMMITTEE INC&DEM&27579&18678797 \tabularnewline
&DNC SERVICES CORPORATION/DEMOCRATIC NATIONAL COMMITTEE&DEM&26444&33411231 \tabularnewline
&NATIONAL REPUBLICAN SENATORIAL COMMITTEE&REP&26074&17532726 \tabularnewline
&DEMOCRATIC SENATORIAL CAMPAIGN COMMITTEE&DEM&4801&10665441 \tabularnewline
&DEMOCRATIC CONGRESSIONAL CAMPAIGN COMMITTEE&DEM&4258&5527039 \tabularnewline
&CLINTON/GORE '96 GEN ELECTION LEGAL \& ACCOUNTING COMPLIANCE&DEM&3643&1926508 \tabularnewline
&DOLE/KEMP '96 COMPLIANCE COMMITTEE INC&REP&3403&2280408 \tabularnewline
2000&BUSH FOR PRESIDENT INC&REP&108472&73718740 \tabularnewline
&REPUBLICAN NATIONAL COMMITTEE - RNC&REP&91994&81798829 \tabularnewline
&NATIONAL REPUBLICAN CONGRESSIONAL COMMITTEE CONTRIBUTIONS&REP&59920&26023012 \tabularnewline
&DNC SERVICES CORPORATION/DEMOCRATIC NATIONAL COMMITTEE&DEM&46366&65189909 \tabularnewline
&GORE 2000 INC&DEM&41427&25422001 \tabularnewline
&NATIONAL REPUBLICAN SENATORIAL COMMITTEE&REP&15283&14461198 \tabularnewline
&GORE/LIEBERMAN GENERAL ELECTION LEGAL AND ACCOUNTING COMPLIANCE FUND&DEM&10384&6410557 \tabularnewline
&BUSH-CHENEY 2000 COMPLIANCE COMMITTEE INC&REP&9539&6770135 \tabularnewline
&DEMOCRATIC CONGRESSIONAL CAMPAIGN COMMITTEE-CONTRIBUTIONS&DEM&6138&8887916 \tabularnewline
&DEMOCRATIC SENATORIAL CAMPAIGN COMMITTEE&DEM&5171&9186161 \tabularnewline
2004&BUSH-CHENEY '04 INC&REP&196987&172686923 \tabularnewline
&JOHN KERRY FOR PRESIDENT, INC&DEM&193913&134195108 \tabularnewline
&REPUBLICAN NATIONAL COMMITTEE&REP&154445&166848971 \tabularnewline
&NATIONAL REPUBLICAN CONGRESSIONAL COMMITTEE&REP&151520&85608871 \tabularnewline
&DNC SERVICES CORPORATION/DEMOCRATIC NATIONAL COMMITTEE&DEM&148660&111769319 \tabularnewline
&KERRY VICTORY 2004&DEM&18892&39900592 \tabularnewline
&MOVEON PAC&DEM&18086&8541254 \tabularnewline
&BUSH-CHENEY '04 COMPLIANCE COMMITTEE INC.&REP&15733&13248329 \tabularnewline
&DEMOCRATIC CONGRESSIONAL CAMPAIGN COMMITTEE&DEM&14445&22942485 \tabularnewline
&KERRY-EDWARDS 2004 INC. GENERAL ELECTION LEGAL AND ACCOUNTING COMPLIANCE FUND&DEM&9449&7849447 \tabularnewline
\bottomrule

\end{tabularx}
\end{table}


\begin{table}[tbp] \centering
\newcolumntype{R}{>{\raggedleft\arraybackslash}X}
\newcolumntype{L}{>{\raggedright\arraybackslash}X}
\newcolumntype{C}{>{\centering\arraybackslash}X}

\ContinuedFloat
\caption[]{Top 10 Selected Committees by Election Year}
\footnotesize
\begin{tabularx}{\linewidth}{c L c r r}

\toprule
{Year}&{Committee Name}&{Party}&{N}&{Amount} \tabularnewline
\midrule \addlinespace[\belowrulesep]
2008&OBAMA FOR AMERICA&DEM&586385&305375905 \tabularnewline
&JOHN MCCAIN 2008 INC.&REP&192852&113007118 \tabularnewline
&REPUBLICAN NATIONAL COMMITTEE&REP&153308&102676723 \tabularnewline
&OBAMA VICTORY FUND&DEM&100727&173782157 \tabularnewline
&MCCAIN-PALIN VICTORY 2008&REP&71786&70354810 \tabularnewline
&NATIONAL REPUBLICAN CONGRESSIONAL COMMITTEE&REP&52022&38212496 \tabularnewline
&DNC SERVICES CORPORATION/DEMOCRATIC NATIONAL COMMITTEE&DEM&47724&43075410 \tabularnewline
&DEMOCRATIC SENATORIAL CAMPAIGN COMMITTEE&DEM&37653&75656317 \tabularnewline
&DEMOCRATIC CONGRESSIONAL CAMPAIGN COMMITTEE&DEM&26887&54122643 \tabularnewline
&NATIONAL REPUBLICAN SENATORIAL COMMITTEE&REP&24163&36781517 \tabularnewline
2012&OBAMA FOR AMERICA&DEM&432226&183162544 \tabularnewline
&ROMNEY FOR PRESIDENT INC.&REP&267105&198944084 \tabularnewline
&OBAMA VICTORY FUND 2012&DEM&199143&337726844 \tabularnewline
&ROMNEY VICTORY INC&REP&153000&427082486 \tabularnewline
&REPUBLICAN NATIONAL COMMITTEE&REP&103368&86067409 \tabularnewline
&DNC SERVICES CORPORATION/DEMOCRATIC NATIONAL COMMITTEE&DEM&47526&32684490 \tabularnewline
&DEMOCRATIC SENATORIAL CAMPAIGN COMMITTEE&DEM&39725&37229191 \tabularnewline
&DEMOCRATIC CONGRESSIONAL CAMPAIGN COMMITTEE&DEM&37588&36622429 \tabularnewline
&NATIONAL REPUBLICAN SENATORIAL COMMITTEE&REP&25625&47530495 \tabularnewline
&NATIONAL REPUBLICAN CONGRESSIONAL COMMITTEE&REP&20237&28755779 \tabularnewline
2016&HILLARY FOR AMERICA&DEM&2470529&275513725 \tabularnewline
&HILLARY VICTORY FUND&DEM&433735&413959093 \tabularnewline
&DCCC&DEM&423193&43381803 \tabularnewline
&REPUBLICAN NATIONAL COMMITTEE&REP&273264&89257279 \tabularnewline
&TRUMP MAKE AMERICA GREAT AGAIN COMMITTEE&REP&245288&67271640 \tabularnewline
&DNC SERVICES CORP./DEM. NAT'L COMMITTEE&DEM&243333&41360116 \tabularnewline
&DSCC&DEM&155483&50151683 \tabularnewline
&DONALD J. TRUMP FOR PRESIDENT, INC.&REP&148068&54199960 \tabularnewline
&NRSC&REP&107949&46881511 \tabularnewline
&MOVEON.ORG POLITICAL ACTION&DEM&62705&6762457 \tabularnewline
\bottomrule

\end{tabularx}
\end{table}


\begin{table}[tbp] \centering
\newcolumntype{R}{>{\raggedleft\arraybackslash}X}
\newcolumntype{L}{>{\raggedright\arraybackslash}X}
\newcolumntype{C}{>{\centering\arraybackslash}X}

\ContinuedFloat
\caption[]{Top 10 Selected Committees by Election Year (Cont.)}

\footnotesize
\begin{tabularx}{\linewidth}{c L c r r}

\toprule
{Year}&{Committee Name}&{Party}&{N}&{Amount} \tabularnewline
\midrule \addlinespace[\belowrulesep]
2020&REPUBLICAN NATIONAL COMMITTEE&REP&1749909&183523920 \tabularnewline
&TRUMP MAKE AMERICA GREAT AGAIN COMMITTEE&REP&1573965&251616570 \tabularnewline
&NRSC&REP&1438060&124814235 \tabularnewline
&WINRED&REP&1212898&-154593632 \tabularnewline
&DCCC&DEM&998346&70370581 \tabularnewline
&DONALD J. TRUMP FOR PRESIDENT, INC.&REP&826138&99219409 \tabularnewline
&DSCC&DEM&764756&128084912 \tabularnewline
&NRCC&REP&508154&49662605 \tabularnewline
&PROGRESSIVE TURNOUT PROJECT&DEM&470252&14016179 \tabularnewline
&BIDEN VICTORY FUND&DEM&428560&474896488 \tabularnewline
\bottomrule

\end{tabularx}
\end{table}


\subsection{Individual Contribution Level Processing}
Once the final set of partisan committees ($\mathcal{M}_{\text{DEM}, t}, \mathcal{M}_{\text{REP}, t}$) is confirmed, individual transactional receipts (\texttt{itcont}) are cleaned and processed according to strict data audit constraints:
\begin{enumerate}
    \item Only explicit, valid individual contribution and refund transaction types are retained (\texttt{10}, \texttt{11}, \texttt{15}, \texttt{15E}, \texttt{21Y}, \texttt{22Y}).
    \item Transactions that are clearly not individual (neither $\texttt{entity\_tp} = \texttt{"IND"}$ nor is missing). Transactions with an \texttt{other\_id} link are dropped to remove simple inter-organizational transfers.
    \item Contributions are filtered by geography to keep only donors with a verified residency code corresponding to the 50 US States or Washington, D.C..
    \item To properly measure total economic support, refund entries (\texttt{21Y}, \texttt{22Y}) are mathematically converted to negative absolute values before aggregation.
\end{enumerate}

\subsection{Aggregation and Variable Metrics}
Valid transactions are aggregated to the state level using the donor's reported state residency code ($s \in \{\text{50 US States}, \text{DC}\}$). For each state $s$ and election cycle $t$, metrics are calculated in two methods ($X$): Net Dollar Amount (\texttt{tot\_amt}, postivie contributions amount net of refunds) and Strictly Positive Contribution Count (\texttt{n\_trans\_pos}, only counting strictly positive contributions). 

The metric for the Republican share of contributions is defined as:
\begin{equation}
    \text{Share}_{X, s, t} = \frac{X_{\text{REP}, s, t}}{X_{\text{REP}, s, t} + X_{\text{DEM}, s, t}}
\end{equation}
To calculate the finalized political alignment variable ($\text{STDIR}_{X, s, t}$), the shares are scaled linearly to fall within the $[-1, 1]$ interval:
\begin{equation}
    \text{STDIR}_{X, s, t} = 2 \times \text{Share}_{X, s, t} - 1
\end{equation}
A final value of $+1$ indicates an entirely Republican contributing base and $-1$ represents an entirely Democratic base.